\begin{solapartat}
\punts{0,2}
$\begin{aligned}[t]
&{}^{222}_{86}Rn \rightarrow {}^{218}_{84}Po + {}^{4}_{2}\alpha\\
&{}^{224}_{86}Rn \rightarrow {}^{224}_{87}Fr + {}^{0}_{-1}\beta^{-1} + {}^{0}_{0}\overline{\nu}_{e^-}
\end{aligned}$

\medskip
$\begin{aligned}[t]
&N(t) = N_0 e^{-\lambda t}\\
&\lambda = \frac{\ln 2}{t_{1/2}}
\end{aligned}$

\medskip
% DUBTE: als subíndexs de les constants de desintegració la pauta escriu «Ra» en lloc de «Rn»; es transcriu tal com és.
\punts{0,2}
$\begin{aligned}[t]
&\lambda_{{}^{222}_{86}Ra} = \frac{\ln 2}{3,82\cdot 24} = 7,56\times 10^{-3}\ \mathrm{h^{-1}}\\
&\lambda_{{}^{224}_{86}Ra} = \frac{\ln 2}{1,80} = 3,85\times 10^{-1}\ \mathrm{h^{-1}}
\end{aligned}$

\medskip
Per al ${}^{222}_{86}Rn$:

% DUBTE: a l'original falta el signe menys: ln(9,00·10^-2/1,00·10^-1) és negatiu i hauria de ser igual a −7,56·10^-3 t; el resultat t = 13,9 h és correcte.
\punts{0,2}
$\begin{aligned}[t]
&9,00\times 10^{-2} = 1,00\times 10^{-1} e^{-7,56\times 10^{-3} t} \Rightarrow\\
&\ln\left(\frac{9,00\times 10^{-2}}{1,00\times 10^{-1}}\right) = 7,56\times 10^{-3} t \Rightarrow t = 13,9\ \mathrm{h}
\end{aligned}$

\medskip
\punts{0,2} Per al ${}^{224}_{86}Rn$: \quad $N(t = 13,9\,\mathrm{h}) = 1,00\times 10^{-1} e^{-3,85\times 10^{-1}\cdot 13,9} = 4,67\times 10^{-4}\ \mathrm{mol}$

\medskip
\punts{0,2} $4,67\times 10^{-4}\ \mathrm{mol} \times \dfrac{6,022\times 10^{23}\ \text{àtoms}}{1\ \mathrm{mol}} = 2,81\times 10^{20}\ \text{àtoms}$
\end{solapartat}

\begin{solapartat}
\punts{0,4} $E = 5,590\ \mathrm{MeV} \times \dfrac{1,60\times 10^{-13}\ \mathrm{J}}{1\ \mathrm{MeV}} = 8,94\times 10^{-13}\ \mathrm{J}$

\medskip
\punts{0,2} $E = \Delta m\cdot c^2$

\medskip
\punts{0,4} $\Delta m = \dfrac{E}{c^2} = \dfrac{8,94\times 10^{-13}}{\left(3,00\times 10^{8}\right)^2} = 9,93\times 10^{-30}\ \mathrm{kg}$
\end{solapartat}
